\section{2026 年 11 月核验增补}

本节按论文首次提交时间排列。每项方法附有可核验的外部入选依据。

% METHOD: 2410.14169
\paragraph{DaRePlane：用方向选择小波平面建模动态形变}
2024.10~\cite{2410_14169}。把时空特征直接分解成几组平面，往往难以表达不同方向的细节。DaRePlane 学习六个方向的双树复小波系数，通过逆变换还原平面并采样；接入高斯时，小型多头网络用这些特征预测时空形变。可学习稀疏掩码约束方向展开引入的额外存储。

原文既评价 NeRF 也评价高斯实现。Plenoptic 和 D-NeRF 的常规动态比较主要展示 NeRF 路线，外科场景另有 DaReNeRF 和 DaReGS 对照。不能用 NeRF 的画质或训练成本直接代替高斯版。作者页注明部分工作在 United Imaging Intelligence 实习期间完成。

\textbf{入选依据} 论文署名单位包括 United Imaging Intelligence（\href{https://arxiv.org/html/2410.14169}{论文原文}）。

% METHOD: 2503.07946
\paragraph{7DGS：把位置、时间和观察方向写进同一高斯}
2025.03~\cite{2503_07946}。7DGS 将三维位置、一维时间和三维观察方向联合表示为七维高斯。在具体时刻和朝向下做条件切片，得到普通光栅化可以使用的三维原语，使几何运动与方向相关外观不必依赖彼此分离的模型。论文还比较附加的自适应高斯细化模块。

原文在 D-NeRF、Technicolor 和六个物理渲染的 7DGS-PBR 场景比较；PBR 特意包含散射和复杂反射等效应。论文总结的超过 400 帧每秒对应关闭自适应细化，不能推广到启用该模块的所有配置。原文署名单位为 United Imaging Intelligence。

\textbf{入选依据} 论文署名单位包括 United Imaging Intelligence（\href{https://arxiv.org/html/2503.07946v1}{论文原文}）。

% METHOD: 2607.18539
\paragraph{AniGS：迭代扩充视角监督，将静态场景转为动态高斯}
2026.07~\cite{2607_18539}。AniGS 从固定的场景高斯出发，先在部分相机位姿渲染，再用视频扩散先验为选定区域生成运动。扩充相机覆盖后，将修复视频重新写入训练集，迭代优化时序形变场。合成时保留原场景的静态内容，不把所有像素交给生成模型重绘。

论文在五个含大片植被、可沿场景移动的户外场景比较代表性动画方法。场景没有真实动态多视图，FVD 的参考由生成分布获得，并辅以用户研究。因此评价针对动画外观和一致性，不能验证不可逆状态变化或昼夜等大范围外观变化。论文署名列有 Waymo 与 Meta 作者。

\textbf{入选依据} 论文署名单位包括 Meta（\href{https://arxiv.org/html/2607.18539v1}{论文原文}）。以原文作者单位为准，修正此前误记的 Waymo。

% METHOD: 2607.21448
\paragraph{GrainGS：隔离规范几何与形变之间的训练梯度}
2026.07~\cite{2607_21448}。GrainGS 首先训练静态规范锚点，然后给锚点产生的各高斯预测时间形变。停止梯度使形变支路的误差信号不会直接把规范参考拉偏。外观另有规范颜色与时间残差两部分，以缓解明暗变化对几何运动的错误解释。

作者在 D-NeRF 和 DG-Mesh 比较图像保真、渲染速度和体积。D-NeRF 的平均 PSNR 相对最强表中基线提升 0.25 dB，结果仍来自所测数据和训练设置。论文作者页列 Qi Tian 同时署名华为消费者业务组。

\textbf{入选依据} 论文署名单位包括 Huawei（\href{https://arxiv.org/html/2607.21448v1}{论文原文}）。

% METHOD: 2608.21828
\paragraph{按时间变化匹配四维高斯的空间滤波尺度}
2026.08~\cite{2608_21828}。动态高斯的运动会改变其跨相机的投影大小，使用固定抗混叠尺度可能在放大时显露空洞，在缩小时出现锯齿。该方法为原语估计时间与焦距深度比的联合分布，在给定时间据条件分布选取采样尺度，再对投影高斯施加平滑。

论文以 SARO-GS 为主要底座，在 Plenoptic Video 上测试低分辨率训练后的一到四倍输出，另用 D-NeRF 做缩小实验。原文指出为每个原语保存时间与焦距深度比的离散密度矩阵，内存随原语数和两个维度的分箱增长，大场景需要另计这项成本。

\textbf{入选依据} 论文署名单位包括 Samsung（\href{https://arxiv.org/abs/2608.21828}{论文原文}）。

% METHOD: 2609.05888
\paragraph{UniFusion：把各相机各时刻的单目深度对齐后再训练动态高斯}
2026.09~\cite{2609_05888}。多个相机或时间的单目预测具有不同尺度与偏移，单帧深度不能直接当成统一的几何监督。UniFusion 用同步相机的稀疏 SfM 参照联合拟合时空深度场，通过低秩时间权重和沿射线预测的深度偏移维持跨帧共享，并将对齐结果用于动态高斯的初始化和优化监督。

论文在 Ego-Exo4D 按时间留出帧，在 EgoHuman 按相机留出视图，两个数据均使用已知相机位姿。深度对齐也有独立于高斯后端的重投影一致性实验。原文展示快速复杂运动仍可受变形后端限制，改善深度先验不等于取消运动模型的表达瓶颈。

\textbf{入选依据} \href{https://github.com/yongzhelyu/unifusion}{公开实现}已找到与论文对应的代码文件。

% METHOD: 2609.24531
\paragraph{Dynamic Thermal Gaussians：用共享几何组织动态 RGB 和热红外}
2026.09~\cite{2609_24531}。该方法将 RGB 与热红外置于同一规范高斯几何，分别用模态属性、时变形变和结构感知路由渲染。模态专有的辐射变化可以由专门原语和属性解释，同时两类观测共享空间与相机参考。

论文在 DynamicRGBT-Scenes 的十二个真实序列中报告热红外 PSNR 平均 32.24 dB。输入依赖 RGB 视频获得的 COLMAP 位姿及与热帧的人工同步，评测的是配对动态新视角，不是以未标定的纯热红外完成重建。

\textbf{入选依据} \href{https://github.com/linlif1869/dtg}{公开实现}已找到与论文对应的代码文件。

% METHOD: 2609.32175
\paragraph{OneFixer：用同一次自回归生成轨迹训练街景视频修复}
2026.09~\cite{2609_32175}。OneFixer 将部署时的自回归历史纳入训练，在同一条生成轨迹上同时做流匹配与逐帧感知监督。它用单步生成修复输入的街景高斯渲染，训练目标兼顾帧质量和连续外观，而输出属于视频修复结果。

论文分别用 Waymo 和首尔驾驶数据训练，报告视频和逐帧指标，并说明细结构仍可能出现问题。不同数据域的模型训练分开进行，单步时延和图像一致性也不能直接解释为三维几何重建精度。

\textbf{入选依据} 论文署名单位包括 42dot、Hyundai Motor Company（\href{https://arxiv.org/abs/2609.32175}{论文原文}）。

% METHOD: 2609.32711
\paragraph{ProDyGS：以单目深度生成静止相机没有采到的代理视角}
2026.09~\cite{2609_32711}。静止或近乎静止的相机没有跨帧视差，难以用不同帧三角化动态表面。ProDyGS 先用单图模型预测每帧深度，把颜色随深度前向投影到相邻虚拟视角，再与真实相机观测一起约束规范高斯和时间形变网络。

在 DyNeRF，系统用单一相机序列训练，并在未见的真实相机视角评价。预测深度与 COLMAP 坐标的尺度平移对齐是该实验的评价设置。代理图像重投影的是同一帧已有观测，不能提供真实遮挡面证据。

\textbf{入选依据} 论文署名单位包括 Sony（\href{https://arxiv.org/abs/2609.32711}{论文原文}）。

% METHOD: 2609.39960
\paragraph{按表示机制理解动态重建的四维综述}
2026.09~\cite{2609_39960}。这篇综述以表示机制而非应用场景组织路线，分别梳理四维 NeRF 的隐式场、特征体和时间先验，以及高斯的显式时空原语、规范形变与逐帧训练。这个框架适合区分运动是直接写入基元，还是经共享规范坐标恢复。

性能章节汇总了动态场景数据、图像质量、几何和效率指标以及论文中的已报告结果。它能说明不同路线的取舍，但这些表格并非统一硬件、输入相机和训练协议下的重新实验。

\textbf{入选依据} 论文署名单位包括 Google DeepMind、Huawei（\href{https://arxiv.org/abs/2609.39960}{论文原文}）。

% METHOD: 2610.01114
\paragraph{仿射对齐图集：在规范高斯构建前吸收组内全局运动}
2026.10~\cite{2610_01114}。作者先用逐帧仿射变换把同一帧组的观测放到扩大的图集坐标，再在对齐后构建规范高斯，最后交由原有时间形变网络拟合剩余变化。这样可减轻跨帧背景位移导致的规范表示模糊，修改集中在初始参考坐标。

论文在 Big Buck Bunny 和 UVG 的十帧组上评估，明显相机运动的片段受益更大，低运动片段改善有限。仿射模型也不能覆盖所有复杂局部形变。Sharp Corporation 署名可直接在作者页核对。

\textbf{入选依据} 论文署名单位包括 Sharp Corporation（\href{https://arxiv.org/html/2610.01114v1}{论文原文}）。

% METHOD: 2610.14685
\paragraph{Demiurge：自回归重渲染时持续合并四维点云记忆}
2026.10~\cite{2610_14685}。Demiurge 把已有单目视频和相机轨迹作为输入，从记忆中投影过去的几何及外观，为新的相机轨迹生成视频，再用补全深度把新观测写回四维点云。高斯重建可以使用这些生成多视图作为下游输入，核心记忆本身是点云。

论文分开评价生成视频的感知质量、跨轨迹同步和几何可重建性。依次执行外观生成、深度恢复与记忆更新，作者报告每条目标轨迹约需五分钟，当前不能视为交互帧率重渲染。

\textbf{入选依据} 论文署名单位包括 Microsoft、Tencent（\href{https://arxiv.org/abs/2610.14685}{论文原文}）。

% METHOD: 2610.24812
\paragraph{SwitchSplat：将物体开关状态从连续时间中分离}
2026.10~\cite{2610_24812}。对可以打开、关闭或切换的物体，时间只是观测顺序，不能单独作为场景状态的索引。SwitchSplat 用共享静态底层储存不变部分，并为每个可交互物体维护可独立切换的二态高斯槽；可微槽表把各帧观测连到相应状态。这使新视图可以组合不同物体的状态，而不必重放某个采集时刻。

实验的常规十个场景与额外两组详细测试区分留出相机和新状态组合。方法当前假定状态为二元，采集约束在单人操作，槽位数量和大致交互窗口还由视觉语言模型给出初值。论文验证的是状态可控渲染，并未建立物理交互模型。

\textbf{入选依据} 论文署名单位包括 Wayve（\href{https://arxiv.org/abs/2610.24812}{论文原文}）。

% METHOD: 2610.35002
\paragraph{DUET4D：用当前高斯约束生成，再筛选伪观测回填}
2026.10~\cite{2610_35002}。DUET4D 先从当前四维高斯抽取跨时空对应，让冻结的单视图扩散模型恢复伪观测时沿这些锚点保持一致。生成结果再反向辅助建图，但只有置信度足够的差异才用于增密，与真实观测持续冲突的高斯则会被裁掉。

作者在三个稀疏视角动态数据上比较渲染效果。初始几何严重失真、运动太快或遮挡时间太长时，对应锚点的支持会下降。系统此时退回较独立的单视图恢复，可以限制错误传播，但无法保证重建真实不可见内容。

\textbf{入选依据} 论文署名单位包括 XGRIDS（\href{https://arxiv.org/abs/2610.35002}{论文原文}）。
